\chapter{By geometry: solids, surfaces \& domains}\label{ch:geometry}
\begin{goals}
\textbf{In this chapter} you classify records against geometry: select or flag what lies inside a wireframe solid, above or below a surface, or between two surfaces; write the fraction of each block a solid fills; and code a whole model by its domain solids in one run.
\end{goals}

A wireframe says where something is. A block model, a set of composites or a drillhole table says what is there. The work of coding a model — which blocks are oxide, which samples are inside the ore shell, how much of each bench block lies below topography — is the work of bringing the two together. Earlier chapters used the results as if they were given — a fraction column such as \col{PARTIAL} in a density, a flag such as \col{IN\_PIT} in a filter (\cref{sec:gt-units,sec:xtab-out}). This chapter is where those columns come from. micro makes them with two dialogs.

\begin{itemize}
\item \menu{Tools>Select by solid / surface…} classifies records against \emph{one} piece of geometry: a closed solid, an open surface, or a gridded surface.
\item \menu{Tools>Assign domains…} classifies a block model against \emph{several} solids at once and writes one domain column.
\end{itemize}

Both are also in the command palette (\key{⌘K}, type \code{solid} or \code{domains}), and the first opens from a layer's own menu with that layer already chosen: \menu{Geometry>Select by this solid…} on a mesh, \menu{Geometry>Flag / ratio by this surface…} on a grid.

\section{What counts as geometry}\label{sec:geometry-kinds}

The \ui{geometry} list at the top of the dialog holds every layer that can classify:

\begin{description}
\item[A mesh held exactly] — a wireframe from \file{.msh}, \file{.obj}, \file{.dxf}, a Datamine pair, a small \file{.ply} (\cref{ch:meshes}). A very large mesh drawn \emph{streamed} keeps no triangles on the page and is not offered; convert it first with \menu{Data \& file>Convert to exact} on its menu.
\item[A grid] — a GeoTIFF, Surfer or ESRI ASCII surface (\cref{ch:grids}), marked ▦ in the list.
\end{description}

When you pick a mesh, micro reads its topology and a line under the list tells you what it found. That line decides the vocabulary of the rest of the dialog:

\begin{table}[htbp]\small
\begin{tabularx}{\linewidth}{@{}LL@{}}\toprule
the dialog says & the choices under \ui{which} \\ \midrule
\emph{closed solid — inside / outside apply} & \ui{inside}, \ui{outside} \\
\emph{open surface (N boundary edges) — above / below / between apply} & \ui{below}, \ui{above}, \ui{between} \\
\emph{gridded surface (nx×ny heightfield) — above / below / between} & \ui{below}, \ui{above}, \ui{between} \\
\emph{forced closed volume (N boundary edges bridged by winding)} & \ui{inside}, \ui{outside} \\ \bottomrule
\end{tabularx}
\caption{What micro reads from the geometry, and the question it lets you ask.}\label{tab:geometry-topology}
\end{table}

A mesh with any boundary edges — edges that belong to one triangle only — is treated as an open surface. That is right for topography and wrong for an ore shell with a few pinholes left by an export. For the second case a \ui{volume} row appears with one checkbox, \ui{treat as closed volume}. Tick it and the mesh is evaluated as a solid; the pinholes are bridged by the winding number, described next.

\begin{tip}
Not sure whether a wireframe is closed? \menu{Analyse>Mesh QA / volume…} on its menu tells you, and steps the camera to each hole (\cref{ch:qa}).
\end{tip}

\section{How “inside” is decided}\label{sec:geometry-winding}

micro decides containment with the generalized \term{winding number} of \textcite{jacobson2013}. For any point, it adds up the solid angle that every triangle of the mesh subtends as seen from that point, signed by which way the triangle faces. From inside a closed solid the triangles surround the point completely and the sum is one full sphere: winding number 1. From outside, the near and far faces cancel: winding number 0. A point is inside when its winding number is above one half.

\begin{figure}[htbp]\centering
\begin{tikzpicture}[font=\sffamily\footnotesize, every node/.style={text=ink}, >=Stealth]
% closed solid
\begin{scope}
\draw[amber, line width=1pt, fill=tint] (0,0.3) .. controls (0.4,1.9) and (2.6,2.2) .. (3.4,1.2) .. controls (4.0,0.3) and (3.0,-0.9) .. (1.8,-0.8) .. controls (0.8,-0.8) and (-0.2,-0.5) .. (0,0.3) -- cycle;
\fill[ok] (1.7,0.6) circle (1.6pt);
\foreach \a in {20,65,110,155,200,245,290,335} \draw[ok, ->, line width=0.4pt] (1.7,0.6) -- ++(\a:0.75);
\node[anchor=north, text=ok] at (1.7,0.05) {$w = 1$};
\fill[fault] (4.9,1.5) circle (1.6pt);
\draw[fault, ->, line width=0.4pt] (4.9,1.5) -- (3.55,1.15);
\draw[fault, ->, line width=0.4pt, dashed] (3.45,1.12) -- (0.25,0.35);
\node[anchor=west, text=fault] at (5.0,1.5) {$w = 0$};
\node[anchor=north, text=muted] at (1.8,-1.1) {a closed solid};
\end{scope}
% holed solid
\begin{scope}[xshift=70mm]
\draw[amber, line width=1pt, fill=tint] (3.0,1.75) .. controls (3.7,1.3) and (3.9,0.2) .. (3.2,-0.5) .. controls (2.6,-1.0) and (0.6,-0.9) .. (0.1,-0.2) .. controls (-0.3,0.5) and (0.3,1.7) .. (1.3,1.95);
\draw[rule, line width=0.6pt, dashed] (1.3,1.95) -- (3.0,1.75);
\fill[ok] (1.9,0.5) circle (1.6pt);
\node[anchor=north, text=ok] at (1.9,0.35) {$w \approx 0.9$};
\node[anchor=south, text=muted] at (2.15,1.95) {a hole};
\node[anchor=north, text=muted] at (1.8,-1.1) {a solid with a gap};
\end{scope}
\end{tikzpicture}
\caption{The winding number in section. Inside a closed solid every direction meets the surface facing outward, and the sum is 1; outside, each ray enters and leaves, and the two faces cancel to 0. A gap removes only the solid angle of the gap, so a point well inside a holed mesh still reads close to 1 and is classified correctly.}\label{fig:geometry-winding}
\end{figure}

This is why the method is robust. A ray-counting test fails the moment a ray passes through a hole or grazes an edge. The winding number degrades gently: a hole takes away only the small solid angle the hole itself subtends, and a solid whose faces are wound inside-out is recognized and handled. Real wireframes — exported, clipped, merged — are rarely perfect, and a classifier that needs them to be perfect answers wrongly without saying so.

Two further choices in the dialog concern how the sum is computed, not what it means.

\begin{table}[htbp]\small
\begin{tabularx}{\linewidth}{@{}lMD@{}}\toprule
setting & meaning & default \\ \midrule
\ui{compute} & \ui{GPU} evaluates on the graphics card. \ui{CPU} runs the same evaluation in a background worker: slower, but it cannot stall a graphics driver on a very large solid. & \ui{GPU} \\
\ui{engine} & \ui{winding} is the robust sum described above; it tolerates holes and flipped faces. \ui{peel · fast} casts one ray down each column of blocks and is exact along Z — much faster for a \emph{clean} solid or surface. Meshes only; a grid has its own method. On a rotated model the rays run along the model's own vertical axis, so peel is used for a closed solid only, and an open surface is evaluated with \ui{winding}. & \ui{winding} \\ \bottomrule
\end{tabularx}
\caption{How the evaluation runs.}\label{tab:geometry-engine}
\end{table}

\begin{caution}
\ui{peel} trusts the mesh. If it reports that columns overflowed, or the result looks wrong on a messy wireframe, switch back to \ui{winding}. The progress line names the engine as it sweeps, and a fraction column's final status line names it again.
\end{caution}

\section{Select, flag, or fraction}\label{sec:geometry-do}

The \ui{do} row chooses what happens to the records that qualify.

\subsection{Select}

\ui{select} sweeps the qualifying records into the ordinary selection, tinted gold in the view (\cref{ch:picking}). Under \ui{layers}, tick the layers to classify — block models, points, drillholes, mesh vertices; the ones currently visible start ticked. \ui{combine} decides how the result meets a selection you already have: \ui{replace}, \ui{add} or \ui{subtract}.

A selection is the lightest result. Use it to look (“which composites are inside the shell?”), to feed an analysis window set to the selection scope (\cref{ch:analysis}), or as the first step of a compound question: select inside the ore shell, then run the dialog again with \ui{below} the weathering surface and \ui{combine} set to \ui{subtract}.

\subsection{Flag column}

\ui{flag column} writes the answer down. It creates a new category column on each ticked layer and puts a value in it for every record that qualifies; the others stay blank.

\begin{table}[htbp]\small
\begin{tabularx}{\linewidth}{@{}lMD@{}}\toprule
setting & meaning & default \\ \midrule
\ui{layers} & The layers that get the column. Only table-backed layers can take one; the others are grayed, with a tooltip. & the visible ones \\
\ui{column} & The column's name. & \code{IN\_}, \code{BELOW\_}, \code{ABOVE\_} or \code{BTW\_} plus the geometry's layer name (see below) \\
\ui{value} & The value written for records that qualify. & \code{IN} \\ \bottomrule
\end{tabularx}
\caption{Settings for a flag column.}\label{tab:geometry-flag}
\end{table}

The offered name is built from the geometry's layer name: its file extension dropped, every run of spaces or punctuation turned into one underscore, the whole in capitals and cut at 24 characters. The demo's \ui{ore shell} therefore offers \code{IN\_ORE\_SHELL}, and a surface named \file{base of ox.dxf} offers \code{BELOW\_BASE\_OF\_OX}. Type a name of your own to override it.

A record is judged by one point: a block by its centroid, a drillhole interval by its position. The result is a painted column like any other (\cref{ch:picking}): the filter sees it (\code{IN\_ORE\_SHELL = "IN"}), the attribute table shows it, export carries it, and you can touch it up by hand with the brush where you disagree with the wireframe. When the run ends the layer is colored by the new column and the status line reports the count.

\subsection{Fraction column}

A centroid is either in or out. For a block cut by a contact that is a crude answer, and the \ui{fraction column} gives a better one: the proportion of each block's volume that qualifies, from 0 to 1.

The block is sampled at a small lattice of probe points, and the fraction is the share of probes that qualify. You set the number of probes per axis in the \ui{probes} row, and the dialog works out what that means in meters.

\begin{table}[htbp]\small
\begin{tabularx}{\linewidth}{@{}lMD@{}}\toprule
setting & meaning & default \\ \midrule
\ui{model} & The block model that gets the column. & the first one \\
\ui{column} & The column's name. & as for a flag \\
\ui{probes} & Probe points along x × y × z in each block. The tally beside it reads, for example, “= 18 probes/block · 3.33×3.33×2.5 m blocklings”. & 3 × 3 × 3 \\ \bottomrule
\end{tabularx}
\caption{Settings for a fraction column.}\label{tab:geometry-ratio}
\end{table}

Probes are per axis because blocks are rarely cubes. A 10 × 10 × 5 m bench block is well served by 3 × 3 × 2; a model that must resolve a thin dyke may want 10 × 10 × 10, at a proportional cost in time. With a closed solid the fraction is always the part \emph{inside}, so the \ui{which} row disappears; with a surface you still choose below, above or between.

The column is numeric, wears a \% badge in the column lists, and is the natural weight for everything downstream. A grade–tonnage series weighted by \code{SG * IN\_ORE\_SHELL} reports the tonnage inside the shell with the edge blocks counted for the part that is in (\cref{ch:gt}).

\begin{note}
Fractions need a block model, axis-aligned or rotated, regular or sub-blocked. On a rotated model micro turns the geometry, not the model: it carries the solid or surface into the model's own lattice frame, where the rotated blocks line up with the axes again, and probes there. The probes are therefore laid out along the block's own edges, and a rotated copy of a model and its solid gives exactly the fractions the unrotated pair gives. A surface's skirt is dropped straight down (or up) in the world before the turn, so \ui{below} keeps meaning below. Select and flag on a rotated model take the same path, judging each block by its centroid as on any model.
\end{note}

\begin{note}
A sub-blocked model has no lattice of whole blocks to sweep, so each block is probed in its own box: the same probe pattern, stretched to the block's own size. A 1~m sub-block and its 20~m parent each get 3~×~3~×~3 probes, and the tally beside \ui{probes} says \emph{in each sub-block's own box}. The work is one pass over the blocks whose box reaches the solid's bounding box; every other block is 0 without a probe. That pass is shared between background workers, one fewer than the computer's processor threads and at most eight: the blocks are cut into slices, each slice is probed whole by whichever worker is free, and the answers are written back in block order, so the column is the same however many workers made it. The status line says how many, as in \ui{(8 workers)}. Peel needs whole columns of blocks, so on a sub-blocked model \ui{peel} is answered by \ui{winding}, block by block, and the tally says so. All of this holds for a rotated sub-blocked model too, such as an Isatis or Datamine export with sub-cells, which used to be refused: a rotated sub-blocked copy of a model gives the same fractions, flags, domains and factor tonnages as its axis-aligned twin. A gridded surface (\cref{sec:geometry-surfaces}) probes each sub-block's own box in the same way.
\end{note}

\begin{tryit}
\begin{enumerate}
\item Open the demo project. Right-click \ui{ore shell} in the layers panel and choose \menu{Geometry>Select by this solid…}. The dialog reports \emph{closed solid — inside / outside apply}.
\item Leave \ui{do} on \ui{select} and \ui{which} on \ui{inside}. Under \ui{layers}, tick \file{deposit.csv} and the drillholes (\ui{holes · demo\_assay.csv}) and untick anything else.
\item Press \ui{Run}. The blocks and intervals inside the shell turn gold — 81,552 blocks and 267 of the 951 intervals.
\item Press \key{Esc} to clear the selection.
\item Open the dialog again and set \ui{do} to \ui{flag column}. Leave the name as offered and the value as \code{IN}, and press \ui{Run}.
\item The status line ends with the expression to use: \code{filter with IN\_ORE\_SHELL == "IN"}. Select \file{deposit.csv} in the layers panel and type that expression in the filter box (micro writes \code{==}, which the filter reads as \code{=}; \cref{app:expr}). Only the blocks inside the shell remain.
\item Open the dialog once more. Set \ui{do} to \ui{fraction column} and \ui{probes} to 3 × 3 × 2, give the column a name of your own such as \code{ORE\_FRAC}, and press \ui{Run}.
\item Color the model by \col{ORE\_FRAC}. The interior is 1, the outside 0, and the skin of the shell shows every value between.
\end{enumerate}
\end{tryit}

\shot{geometry-dialog}{Select by solid on the demo's ore shell, set to write a fraction column, \col{ORE\_FRAC}, onto \file{deposit.csv} with 3 × 3 × 2 probes per block — on the demo's 10 × 10 × 5 m blocks, “18 probes/block · 3.33×3.33×2.5 m blocklings”.}

\section{Surfaces: above, below, between}\label{sec:geometry-surfaces}

An open surface — topography, the base of weathering, a fault — has no inside. It divides space into two sides, and the dialog's vocabulary changes to match: \ui{below}, \ui{above}, and \ui{between}.

To reuse the winding number, micro closes the surface. From every boundary edge it drops a vertical wall, a \emph{skirt}, straight down (for \ui{below}) or straight up (for \ui{above}) to a flat cap far away, which turns “below the surface” into “inside a closed body”. Because the walls are vertical, the closed body has the surface's own footprint at every depth, and the same robust containment applies however far below the surface a record lies.

\begin{figure}[htbp]\centering
\begin{tikzpicture}[font=\sffamily\footnotesize, every node/.style={text=ink}]
\fill[tint] (0.4,2.1) .. controls (1.6,2.7) and (2.6,1.5) .. (3.8,2.0) .. controls (4.8,2.4) and (5.6,1.9) .. (6.4,2.2) -- (6.4,-1.4) -- (0.4,-1.4) -- cycle;
\draw[rule, dashed, line width=0.5pt] (0.4,2.1) -- (0.4,-1.4) -- (6.4,-1.4) -- (6.4,2.2);
\draw[amber, line width=1.1pt] (0.4,2.1) .. controls (1.6,2.7) and (2.6,1.5) .. (3.8,2.0) .. controls (4.8,2.4) and (5.6,1.9) .. (6.4,2.2);
\node[anchor=south, text=amber] at (3.4,2.35) {the surface};
\node[text=muted] at (3.4,-1.75) {the cap, far below};
\node[anchor=east, text=muted] at (6.3,0.2) {the skirt};
\fill[ok] (3.0,1.0) circle (1.6pt); \node[anchor=west, text=ok] at (3.1,1.0) {below};
\fill[info] (2.2,2.9) circle (1.6pt); \node[anchor=west, text=info] at (2.3,2.9) {above};
\fill[faint] (7.3,1.2) circle (1.6pt); \node[anchor=west, text=muted] at (7.4,1.2) {beyond the footprint: unclassified};
\end{tikzpicture}
\caption{An open surface closed for \ui{below}: vertical walls from its boundary down to a flat cap far below. A point under the surface is inside the closed body. A point outside the surface's footprint is on neither side, and is left unclassified.}\label{fig:geometry-skirt}
\end{figure}

A record that lies beyond the surface's footprint — east of where the topography was surveyed, say — is neither above nor below it. micro leaves such records unclassified instead of extending the surface by guesswork. If your surface is smaller than your model, extend it first: \menu{Geometry>Extend to footprint…} on the surface's menu (\cref{ch:grids}).

\ui{between} needs a second surface. Choosing it adds a \ui{lower} row listing the other open surfaces, and a note stating the meaning precisely: \emph{between = below the first surface and above this one}. Pick topography as the geometry and the base of oxidation as \ui{lower}, and the records flagged are the oxide zone. If no second surface is loaded, the dialog says so and \ui{Run} stays disabled.

\subsection{Gridded surfaces}

A grid is a surface with a regular footprint, and micro uses that. There is no mesh to wind around: for each record it samples the grid's elevation at the record's x and y, bilinearly between the four surrounding nodes, and compares it with the record's z. For a fraction column it compares the surface with the block's vertical extent at each probe position. A block of a rotated model has no vertical extent to compare, so there the probes fill the block in three dimensions: they are laid out along the block's own edges, and each one is compared with the surface at its own x and y. The grid itself is never turned. A probe count of 1 along z, the default on an axis-aligned model, becomes the larger of the x and y counts on a rotated one, and select and flag use 2 × 2 × 2 probes and a majority. The dialog says \emph{a bilinear sample, no winding}, and the \ui{engine} row disappears.

Two details matter for trust:

\begin{itemize}
\item A large GeoTIFF is displayed through a reduced overview, but it is \emph{evaluated} at full resolution. Before the run micro reads the full-resolution window of the grid that covers the target layers, and the status line says so. If that read fails it falls back to the displayed grid and tells you.
\item \ui{between} with a grid needs a second grid; a grid and a mesh cannot be mixed in one \ui{between}.
\end{itemize}

\begin{tryit}
\begin{enumerate}
\item In the demo project choose \menu{Tools>Select by solid / surface…} and set \ui{geometry} to \ui{topography}.
\item The choices are now \ui{below}, \ui{above}, \ui{between}. Choose \ui{flag column} and \ui{above}, name the column \code{AIR} and set the value to \code{AIR}.
\item Tick \file{deposit.csv} only and press \ui{Run}.
\item Filter with \code{AIR = "AIR"} to see which block centroids sit above the ground. The demo's model has a flat top bench while its synthetic topography rolls by several meters, so a scatter of top-bench blocks shows up — 1,996 of them. In a model you are handed, every one of them is a question for whoever built it — the check is worth running.
\end{enumerate}
The topography appears in the list as \ui{▦ topography}: although it is drawn as shaded relief, it is a grid, so it is evaluated by bilinear sampling.
\end{tryit}

\section{Assign domains}\label{sec:geometry-domains}

A\index{domain!assigning} coded model has one column saying which domain each block belongs to. Building it by hand means one flag per solid and then a rule for the blocks that two solids both claim. \menu{Tools>Assign domains…} does the whole loop in one run and applies one clear rule: \emph{each block goes to the solid holding its largest fraction}.

\begin{table}[htbp]\small
\begin{tabularx}{\linewidth}{@{}lMD@{}}\toprule
setting & meaning & default \\ \midrule
\ui{solids} & A checkbox per exact mesh. An open surface is listed but grayed — “open surface, claims nothing”. & all ticked \\
\ui{model} & The regular block model to code. & the first one \\
\ui{column} & The name of the domain column. & \code{DOMAIN} \\
\ui{probes} & Probe points per block along x × y × z, as for a fraction column. & 2 × 2 × 2 \\ \bottomrule
\end{tabularx}
\caption{The assign-domains dialog.}\label{tab:geometry-domains}
\end{table}

Press \ui{Assign}. micro evaluates the fraction of every block inside every ticked solid, then gives each block the name of the solid with the largest share. A block on a contact lands with its dominant domain. A block no solid touches stays unpainted. The values of the column are the solids' layer names without their file extension, so name your layers the way you want your domains to read.

The result is an ordinary painted column: filter by it, color by it, use it as the group in a resource cross-tab (\cref{ch:resources}) or as the split in a composite (\cref{ch:drillholes}), and correct individual blocks with the brush. The status line reports how many blocks were claimed, out of how many.

\begin{caution}
The rule is “largest fraction”, not “first in the list” and not “any overlap”. Where two solids overlap, the overlap goes to whichever fills more of the block; it is never counted twice. If your domains are meant to nest — a high-grade core inside a wider envelope — the core does not automatically win a block that the envelope also fills completely: on an exact tie the solid listed first takes the block. Flag the core separately, or cut it out of the envelope upstream.
\end{caution}

\begin{note}
Assign domains needs a block model. A rotated or sub-blocked one works as for a fraction column.
\end{note}

\section{Recipes, and in a routine}\label{sec:geometry-routine}

The select-by-solid dialog has a \ui{save as recipe…} button beside \ui{Run}. It stores the whole evaluation — geometry, action, side, engine, targets, column and value — as a recipe that names its layers, so “flag the oxide onto this month's model” becomes a file that runs again from \menu{Tools>Recipes} (\cref{ch:recipes}). Opening such a recipe for preview shows the dialog configured and not yet run.

In a routine the same evaluation is the \code{solid} tool. A fraction step's \code{model} is checked before anything runs; a geometry, lower surface or target layer that is not loaded stops the step itself, with a message naming the layer. Name a layer by its file, as \code{ore\_shell.obj} below, or by its label in the layers panel (\ui{ore shell}); either finds it.

\begin{lstlisting}[language=yaml]
- tool: "solid"
  with:
    geometry: "ore_shell.obj"
    do: "flag"                 # select | flag | ratio
    which: "inside"
    targets:
      - "deposit.csv"
    column: "ORE"
    value: "IN"
- tool: "solid"
  with:
    geometry: "ore_shell.obj"
    do: "ratio"
    model: "deposit.csv"
    column: "ORE_FRAC"
    probes:
      - 3
      - 3
      - 2
\end{lstlisting}

\begin{table}[htbp]\small
\begin{tabularx}{\linewidth}{@{}lMD@{}}\toprule
setting & meaning & default \\ \midrule
\code{geometry} & the solid, surface or grid, by layer name & required \\
\code{do} & \code{select}, \code{flag} or \code{ratio} & \code{flag} \\
\code{which} & \code{inside} / \code{outside}, or \code{below} / \code{above} / \code{between} & the first that applies \\
\code{lower} & the second surface, for \code{between} & — \\
\code{targets} & the layers to classify (select, flag) & the visible ones \\
\code{combine} & \code{replace}, \code{add} or \code{subtract} (select) & \code{replace} \\
\code{column} & the column to write (flag, ratio) & as the dialog offers \\
\code{value} & the value written (flag) & \code{IN} \\
\code{model} & the block model (ratio) & required for ratio \\
\code{probes} & probes per axis (ratio): a list of three, or one number \code{n} for n × n × n. Anything else stops the step with an error & 3, 3, 3 \\
\code{engine} & \code{peel} for the fast engine & winding \\
\code{forceClosed} & \code{true} to treat a holed mesh as a closed volume & as the layer is \\ \bottomrule
\end{tabularx}
\caption{Settings of the \code{solid} tool. They are the dialog's own, by name. The step runs through the dialog itself rather than the tool registry, so it is not in the generated reference of \cref{app:tools}; this table is its reference.}\label{tab:geometry-tool}
\end{table}

A flag or ratio step lists the column it writes in the routine's preview — with “(replaces)” when the column already exists — so the opt-in no-overwrites guard can stop a run that would write over last month's flags (\cref{ch:routines}). Combine it with a check on a later analysis step: flag the shell, run grade–tonnage filtered to it, and require that the tonnage has not moved by more than a tolerance.

Assign domains has no recipe of its own. To automate a coded model, use one \code{ratio} step per solid, then a calculated column that picks the largest.
